%% toc: P2 21 21.1,21.8
\section{Proofs for validity, finite satisfiability and semi-conservative reductions}
\label{sec:annex-p02-w02b-p02-031-validity-finsat}

Sets of formulas are sets of codes under the fixed effective coding of
\cref{sec:decision-hierarchy-01}, and the classes \(\mathrm{Val}\),
\(\mathrm{Fin\text{-}sat}\) and \(\mathrm{Non\text{-}sat}\) are those of
Definition~\ref{def:p02-w02-p02-036-classes}.  Under that coding the set
\(\mathrm{FO}\) of codes of formulas, and for each finite signature the set of
codes of its sentences, are decidable.

\begin{annexitem}{annex:p02-w02b-p02-031-val-to-finsat}{thm:p02-w02b-p02-031-val-to-finsat}{Validity reduces to finite satisfiability}
\begin{proof}
Let \(R_{\Lang}\) be a binary relation symbol of \(\Lang\).

\emph{Step 1: universal closure.}  For a formula \(\varphi\) whose free
variables, listed in their order of first occurrence, are
\(x_1,\ldots,x_p\), put \(\varphi^{\forall}\defeq\forall x_1\cdots\forall
x_p\,\varphi\), a sentence.  The free variables are computed by recursion on
\(\varphi\), so \(\varphi\mapsto\varphi^{\forall}\) is computable.  By the
satisfaction clause for \(\forall\), applied \(p\) times, and
\cref{lem:fol-coincidence}, a structure \(\mathfrak A\) satisfies
\(\varphi^{\forall}\) iff \(\mathfrak A,s\entails\varphi\) for every
assignment \(s\).  Hence \(\varphi\) is valid iff \(\varphi^{\forall}\) is
true in every structure for the signature of \(\varphi\).  By
\cref{lem:reduct-invariance} this does not change when larger signatures are
admitted: a structure for a larger signature satisfies \(\varphi^{\forall}\)
iff its reduct does, and every structure for the signature of \(\varphi\) has
expansions.  So \(\varphi\in\mathrm{Val}\) iff \(\varphi^{\forall}\) belongs
to the set \(V\) of codes of valid sentences.

\emph{Step 2: a recognizer.}  The set \(V\) is computably enumerable
(\cref{thm:fol-validity-enumerable}), so some Turing machine \(\mathcal R\)
recognizes it (\cref{thm:ce-equivalent-forms}): its input alphabet contains
the digits \(0,1\), since it is run on binary numerals, and its computation
on \(\operatorname{bin}(n)\) reaches \(q_{\mathrm{acc}}\) iff \(n\in V\).

\emph{Step 3: the map.}  For \(\varphi\in\mathrm{FO}\) put
\[
 F(\varphi)\defeq
 \bigl(\psi_{\mathcal R,\operatorname{bin}(\ulcorner\varphi^{\forall}\urcorner)}
 \bigr)^{\sharp}[R_{\Lang}/E],
\]
where \(\ulcorner\varphi^{\forall}\urcorner\) is the code of
\(\varphi^{\forall}\).  The run signature \(\Lang_{\mathcal R}\) of
Definition~\ref{def:p02-w02-p02-036-run-sentence} consists of binary and unary
relation symbols, so the path translation of
Definition~\ref{def:p02-w02-p02-036-gadget} applies to its sentences, and
\(F(\varphi)\) is an \(\{R_{\Lang}\}\)-sentence, hence a member of
\(\operatorname{Sent}(\Lang)\).  The map \(F\) is computable as a composite
of computable maps: \(\varphi\mapsto\varphi^{\forall}\) by Step~1; the coding
and \(n\mapsto\operatorname{bin}(n)\); \(w\mapsto\psi_{\mathcal R,w}\) for the
fixed machine \(\mathcal R\) by Lemma~\ref{thm:p02-w02-p02-036-run-sentence};
\(\chi\mapsto\chi^{\sharp}\) by
Lemma~\ref{thm:p02-w02-p02-036-gadget}(3); and the letter-for-letter renaming.

\emph{Step 4: the equivalence.}  For \(\varphi\in\mathrm{FO}\), write
\(w_{\varphi}\defeq\operatorname{bin}(\ulcorner\varphi^{\forall}\urcorner)\).
Then
\[
\begin{aligned}
 \varphi\in\mathrm{Val}
 &\iff \ulcorner\varphi^{\forall}\urcorner\in V
   &&\text{(Step 1)}\\
 &\iff \text{the computation of }\mathcal R\text{ on }w_{\varphi}
   \text{ reaches }q_{\mathrm{acc}}
   &&\text{(Step 2)}\\
 &\iff \psi_{\mathcal R,w_{\varphi}}\in
   \mathrm{Fin\text{-}sat}(\operatorname{Sent}(\Lang_{\mathcal R}))
   &&\text{(Lemma~\ref{thm:p02-w02-p02-036-run-sentence})}\\
 &\iff \psi_{\mathcal R,w_{\varphi}}^{\sharp}\in
   \mathrm{Fin\text{-}sat}(\operatorname{Sent}(\{E\}))
   &&\text{(Lemma~\ref{thm:p02-w02-p02-036-gadget}(3))}\\
 &\iff F(\varphi)\in\mathrm{Fin\text{-}sat}(\operatorname{Sent}(\Lang))
   &&\text{(Lemma~\ref{thm:p02-w02-p02-036-transfer}).}
\end{aligned}
\]

\emph{Step 5: the consequence.}  Suppose a machine decides
\(\mathrm{Fin\text{-}sat}(\operatorname{Sent}(\Lang))\).  Given a number
\(n\), reject if \(n\notin\mathrm{FO}\); otherwise compute \(F(n)\) and give
the decider's answer on it.  By Step~4 this procedure halts on every input and
accepts exactly the codes of valid formulas.
\end{proof}
\end{annexitem}

\begin{annexitem}{annex:p02-w02b-p02-031-semi-inseparable}{thm:p02-w02b-p02-031-semi-inseparable}{Semi-conservative reductions preserve inseparability}
\begin{proof}
Suppose, for a contradiction, that a decidable set \(W\) separates
\(\mathrm{Non\text{-}sat}(Y)\) from \(\mathrm{Fin\text{-}sat}(Y)\), that is,
\(\mathrm{Non\text{-}sat}(Y)\subseteq W\) and
\(W\cap\mathrm{Fin\text{-}sat}(Y)=\varnothing\).  Put
\(W'\defeq\{n\in X:f(n)\in W\}\).  The set \(W'\) is decidable: given \(n\),
decide whether \(n\in X\), which is possible because \(X\) is recursive;
if so, compute \(f(n)\), which halts because \(f\) halts on every member of
\(X\); then decide whether \(f(n)\in W\).

If \(\psi\in\mathrm{Non\text{-}sat}(X)\), then
\(f(\psi)\in\mathrm{Non\text{-}sat}(Y)\subseteq W\) by clause~1 of
Definition~\ref{def:p02-w02b-p02-031-semi-conservative}, so \(\psi\in W'\).
If \(\psi\in\mathrm{Fin\text{-}sat}(X)\), then
\(f(\psi)\in\mathrm{Fin\text{-}sat}(Y)\) by clause~2, and this set is disjoint
from \(W\), so \(f(\psi)\notin W\) and \(\psi\notin W'\).  Hence \(W'\)
separates \(\mathrm{Non\text{-}sat}(X)\) from \(\mathrm{Fin\text{-}sat}(X)\),
against the hypothesis.  So no decidable set separates
\(\mathrm{Non\text{-}sat}(Y)\) from \(\mathrm{Fin\text{-}sat}(Y)\).  By the
symmetry remark after Definition~\ref{def:p02-w02-p02-036-inseparable}, no
decidable set separates them in the other order either, and the two sets are
recursively inseparable.
\end{proof}
\end{annexitem}

\begin{annexitem}{annex:p02-w02b-p02-031-semi-not-conservative}{thm:p02-w02b-p02-031-semi-not-conservative}{A semi-conservative reduction without a conservative one}
\begin{proof}
Each of \(X\) and \(Y\) is the set of one fixed code, so membership is decided
by comparing a number with that code; and the constant map \(f\), which sends
the code of \(\prec_K\) to the code of \(\exists x\,Kxx\), is computable.

By Lemma~\ref{thm:p02-w02-p02-036-infinity-axiom}, \(\prec_K\) has a model and
has no finite model.  Hence \(\mathrm{Non\text{-}sat}(X)=\varnothing\) and
\(\mathrm{Fin\text{-}sat}(X)=\varnothing\), and both clauses of
Definition~\ref{def:p02-w02b-p02-031-semi-conservative} hold for \(f\), since
the image of the empty set is empty.  So \(f\) is a semi-conservative
reduction from \(X\) to \(Y\).

Let \(g:X\to Y\) be any map.  Its only possible value is \(\exists x\,Kxx\),
so \(g(\prec_K)=\exists x\,Kxx\).  The one-element structure with domain
\(\{0\}\) and \(K=\{(0,0)\}\) satisfies \(\exists x\,Kxx\), so \(g(\prec_K)\)
is finitely satisfiable, whereas \(\prec_K\) is not.  The required
biconditional for finite satisfiability therefore fails at
\(\psi=\prec_K\), and no such \(g\) exists.
\end{proof}
\end{annexitem}
